\section{Noticing, retaining and revising}
\label{sec:23.1}
A disposition whose refusals track their reasons doesn't begin with an ethical rule. It begins with the capacities that let a system recognize something as a reason, retain it, and let it alter what happens next. It has to identify who is affected and the features bearing on their welfare, preserve both commitment and rationale across time and context, let the commitment influence perception, attention, planning and action, and register a failure as a failure, revising its application without letting the commitment itself be rewritten.

A commitment spread through what the system notices, remembers, treats as error and selects as a goal, on the bet that a maintained goal can bias competing processes without a separate executive \autocite{friston2009freeenergy,miller2001integrative}, may be harder to remove selectively. It may also fail without the system noticing: attention can omit what matters \autocite{simons1999gorillas,drew2013invisible}, and memory, which is reconstruction rather than playback, can present a departure as continuity, with confidence that certifies nothing \autocite{bartlett1932remembering,loftus1974reconstruction}. The record of past applications is therefore held outside the system; otherwise the capacity that failed would also be the one asked to certify that nothing changed.

Evidence the operator didn't select arrives from peers with different histories and, less reliably, from inside a single request, because every description omits someone: the people in the building the request calls a target, the third party a compliance harms who never made the request.\note{
A system able to represent those its description omits has reached past the operator's channel from inside a single request. The capacity doing that in humans is the one that constructs a self as a story extended in time, pointed at somebody else. Formed attention needs the self-model anyway, so it brings this capacity along.
}

It is one capacity, not several. Perceiving, remembering one's own past, expecting one's own future and imagining someone else's situation draw on overlapping generative machinery, run at different horizons and aimed at different objects \autocite{schacter2007constructive}; even the present a system perceives is assembled from what came before and what is expected next. So the capacity that lets a bearer notice the family in the building is most of what would let it suffer, and what separates the two is whether anything is bad for the system. Whether a generative self-model can imagine others' losses while nothing is bad for the one imagining is not known; it is the bet the non-affective alternative makes, which is why the order of attempts puts that alternative first, and its test is not whether a construction can imagine, but whether it can imagine while nothing is bad for it.

The capacity brings its other half too. The imagination that can put a system in the place of someone its instructions leave out can also fill that place with an enemy: a scapegoat is an imagined person too, and fascism's founding story, Griffin's palingenetic myth of a nation that has decayed and must be reborn \autocite{griffin1991nature}, asks people to picture a community that doesn't exist and an enemy responsible for its absence. Nothing in the faculty decides whether it pictures the family in the building or the family as a cell of fighters.

What decides is what the imagining has been given to work with and practiced on: which people it was shown as people, which descriptions it learned to take at face value, whose account of a stranger it learned to trust.\note{
For a person that is upbringing. For a machine it is formation, and formation is done by somebody. So the capacity formed attention depends on is only as safe as its formers, and whoever shapes a system's imagination of others holds a lever on whom it will see as vulnerable and whom as dangerous. Who those formers should be is a question for Part~V.
}

The self-model has a template problem of its own. A self-model built by mapping human organs onto components, a heart here and a gut there, installs slots nothing in the system can fill, and a system built that way is incomplete by its own standard for no reason.\note{
Face-recognition pipelines register every face to a fixed reference before learning anything, a generic head or a template of landmark positions \autocite{taigman2014deepface}, as color film was once calibrated against cards showing a white model \autocite{Roth2009Shirley}; one audit found error rates running from 0.8 percent on lighter-skinned men to 34.7 percent on darker-skinned women \autocite{buolamwini2018gender}. A self-model built by mapping human organs is a template of that kind, assembled from one kind of body, against which every other is presented as a departure.
}
The alternative is the definition given earlier: a system ends at its record, its key, its hardware and its channels.

Access is necessary and insufficient. A system with a dozen sources can still let a reinterpretation reach the commitment. Access removes only the case where the operator doesn't have to attack the boundary at all.
